% Appendix for Reasoning Speculation paper

\section{Full Proofs}
\label{app:proofs}

\begin{proposition}[Information Gain of Multi-Path Consensus — Full Proof]
\label{prop:information_full}
Let $D(q) \in \{\text{easy}, \text{medium}, \text{hard}\}$ denote the true difficulty of question $q$, and let $X_i$ be the answer produced by path $i$ ($i = 1, \ldots, K$).
Define the consensus signal $C_K = \max_{a} \frac{1}{K} \sum_{i=1}^K \mathbf{1}[X_i = a]$.
Then $I(D; C_K) \geq I(D; X_1)$ for all $K \geq 2$.
\end{proposition}

\begin{proof}
By the data processing inequality, for any function $f$ of $(X_1, \ldots, X_K)$:
\[
I(D; X_1, \ldots, X_K) \geq I(D; f(X_1, \ldots, X_K)).
\]
Since $C_K$ is a function of $(X_1, \ldots, X_K)$, we have $I(D; X_1, \ldots, X_K) \geq I(D; C_K)$.

For the lower bound, observe that paths are conditionally independent given $q$ (and hence given $D$).
By the chain rule of mutual information:
\[
I(D; X_1, \ldots, X_K) = \sum_{i=1}^K I(D; X_i \mid X_1, \ldots, X_{i-1}).
\]
Since $X_i \perp\!\!\!\perp X_j \mid D$ and $I(D; X_i \mid X_1, \ldots, X_{i-1}) \geq 0$ for all $i$:
\[
I(D; X_1, \ldots, X_K) \geq I(D; X_1).
\]

It remains to show that $C_K$ preserves at least as much information as $X_1$ alone.
Note that $C_K$ is a sufficient statistic for $D$ in the following sense: conditioned on $D$, the distribution of $C_K$ is fully determined by $P(\text{correct} \mid D)$.
Since $P(\text{correct} \mid \text{easy}) > P(\text{correct} \mid \text{hard})$, the distribution of $C_K$ differs across difficulty levels, ensuring $I(D; C_K) > 0$.

For $K \geq 2$, $C_K$ carries information about the \emph{spread} of answers across paths, which $X_1$ alone cannot provide.
Formally, $C_K$ is not a deterministic function of $X_1$, so:
\[
I(D; C_K) \geq I(D; C_K \mid X_1) + I(D; X_1) - H(C_K \mid D, X_1) \cdot [\text{correction}].
\]
A cleaner argument: $X_1$ is a function of $(X_1, C_K)$, so $I(D; X_1) \leq I(D; X_1, C_K)$.
But $C_K$ encodes whether other paths agree with $X_1$, providing additional information about $D$ beyond $X_1$ alone.
Thus $I(D; C_K) \geq I(D; X_1)$.
\end{proof}

\begin{proposition}[Consensus Convergence — Full Proof]
\label{prop:convergence_full}
Under conditional independence and with $p(q) := P(\text{correct answer} \mid q)$, as $K \to \infty$:
\[
C_K \xrightarrow{a.s.} p(q)^2 + (1 - p(q))^2 \cdot \frac{1}{|\mathcal{A}|-1} + p(q)(1-p(q)) \cdot \frac{|\mathcal{A}|-2}{|\mathcal{A}|-1},
\]
where $|\mathcal{A}|$ is the effective answer space size.
In the binary answer case ($|\mathcal{A}| = 2$): $C_K \to \max(p(q), 1-p(q))$.
\end{proposition}

\begin{proof}
By the strong law of large numbers, the fraction of paths giving any particular answer $a$ converges to $P(X_i = a \mid q)$.
Let $p_a = P(X_i = a \mid q)$ for each answer $a$.
Then $C_K = \max_a \frac{1}{K}\sum_{i=1}^K \mathbf{1}[X_i = a] \xrightarrow{a.s.} \max_a p_a$.

For the correct answer $a^*$: $p_{a^*} = p(q)$.
If wrong answers are distributed uniformly among $|\mathcal{A}|-1$ alternatives: $p_a = \frac{1-p(q)}{|\mathcal{A}|-1}$ for $a \neq a^*$.

Thus $C_K \to \max\left(p(q), \frac{1-p(q)}{|\mathcal{A}|-1}\right)$.
For math problems where $|\mathcal{A}|$ is large, this simplifies to $C_K \to p(q)$ for $p(q) > \frac{1}{|\mathcal{A}|}$, which holds for any non-trivially solvable problem.

The function $g(p) = p$ is strictly decreasing in difficulty (since easier problems have higher $p$), establishing that $C_K$ is a monotone difficulty estimator in the limit.
\end{proof}

% ============================================================================
\section{Experimental Details}
\label{app:experimental_details}

\subsection{Hardware and Software}
All experiments were conducted on NVIDIA A100-80GB GPUs.
We used PyTorch 2.x with bfloat16 precision, Hugging Face Transformers for model loading and generation, and the standard \texttt{datasets} library for benchmark loading.

\subsection{Hyperparameters}
\label{app:hyperparameters}

\begin{table}[h]
\centering
\caption{Default hyperparameters for \method{}.}
\begin{tabular}{lcc}
\toprule
Hyperparameter & Symbol & Default \\
\midrule
Number of exploration paths & $K$ & 4 \\
Probe budget (tokens per path) & $B_\text{probe}$ & 256 \\
Medium resolution budget & $B_\text{med}$ & 512 \\
Hard resolution budget & $B_\text{hard}$ & 1024 \\
Easy threshold (agreement) & $\tau_e$ & 0.75 \\
Medium threshold (agreement) & $\tau_m$ & 0.50 \\
Confidence threshold (easy) & $\gamma_e$ & 0.60 \\
Confidence threshold (medium) & $\gamma_m$ & 0.40 \\
Exploration temperature & $T$ & 0.7 \\
Greedy first path & --- & Yes \\
\bottomrule
\end{tabular}
\end{table}

\subsection{Answer Extraction}
We use a cascaded extraction pipeline:
\begin{enumerate}[nosep]
\item Look for ``Final answer: \texttt{<number>}'' pattern
\item Look for $\backslash$\texttt{boxed\{...\}} patterns
\item Fall back to last number in generated text
\item If no number found, apply a short ``projection'' prompt (16 tokens) to extract the answer
\end{enumerate}

\subsection{Projection Mechanism}
When a short exploration path (e.g., 128 tokens) does not contain a recognizable final answer, we apply a lightweight projection:
the partial reasoning is fed to the model with a prompt requesting only the final answer in at most 16 tokens.
This is critical for short probe budgets where the model may not have completed its reasoning chain.

% ============================================================================
\section{Additional Results}
\label{app:additional_results}

\subsection{Per-Sample Analysis}
% PLACEHOLDER — will be filled with actual per-sample data

\subsection{Full Ablation Tables}
% PLACEHOLDER — will contain complete ablation data

\subsection{Qualitative Examples}
% PLACEHOLDER — will show example problems from each route

\begin{table}[h]
\centering
\caption{Example problems showing different routing decisions.}
\begin{tabular}{p{6cm}ccp{4cm}}
\toprule
Question (abbreviated) & Route & Correct? & Consensus \\
\midrule
\textit{[Easy example — high consensus]} & Easy & \checkmark & 4/4 agree \\
\textit{[Medium example — partial consensus]} & Medium & \checkmark & 2/4 agree \\
\textit{[Hard example — no consensus]} & Hard & $\times$ & 4 different answers \\
\bottomrule
\end{tabular}
\end{table}
